
\subsubsection{6.1 What the Results Tell Us About Gap Signal in Weight Space}

Three findings emerge robustly from our experiments:

\textbf{Finding 1: Generalization gap is learnable.} Both FlatMLP (r = 0.557) and DWSNet (r = 0.510)
exceed the existence threshold. The A1 audit (Spearman(gap, $-test_acc) = -0.142)$ confirms that
gap carries independent predictive signal. Weight tensors encode the overfitting signature --- the
degree to which a model's training adaptations have drifted toward memorization --- in a form that
learned encoders can extract.

\textbf{Finding 2: Architecture specificity matters, specifically cross-layer reasoning.} NFT's cross-layer
attention is the only encoder that improves over FlatMLP on gap prediction, with directional
statistical confidence. DWSNet and GNN both underperform FlatMLP on gap despite their equivariant
architectures. This dissociation --- NFT better than FlatMLP, DWS/GNN worse --- suggests that the
relevant architectural property is not equivariance per se but the capacity to integrate evidence
across layer boundaries. Generalization gap is a globally distributed property of the weight tensor;
detecting it may require an encoder that can reason about how representations change across consecutive
layers, not just within each layer independently.

\textbf{Finding 3: NFT captures gap-specific overfitting structure (P3, r = 0.73).} NFT's gap predictions
contain information about true generalization gap that is statistically independent of test accuracy
rank. We hypothesize (not yet verified) that this corresponds to inter-layer weight co-variation
patterns that emerge during memorization --- patterns accessible to NFT's cross-layer attention but
filtered out by DWS's within-layer averaging and GNN's fixed graph connectivity. This mechanistic
step was not directly tested in the current work.

\subsubsection{6.2 The FlatMLP Test Accuracy Anomaly}

FlatMLP predicts test accuracy at r = 0.279 in our zoo, substantially below the ~0.85 reported by
Unterthiner et al. [2020] on their zoo. This requires explanation, as it affects the $\Delta$ computation
(RQ4) and the interpretation of the dual-target comparison.

We identify two plausible explanations:

\textbf{Explanation 1 (High plausibility): Zoo generation artifact.} Our zoo was generated with a restricted
hyperparameter grid and a specific CNN architecture variant. Models may cluster at near-saturation
training accuracy, compressing the test\_acc variance and making test\_acc harder to predict from weights.
Gap (which reflects the margin between training and test performance) may retain more variance because
it captures subtle memorization differences among high-accuracy models.

\textbf{Explanation 2 (Medium plausibility): Optimizer sensitivity.} The 3-trial search budget may have
failed to find a good learning rate for FlatMLP's test\_acc regression task specifically, while the
gap regression task is more robust to learning rate choice.

\textbf{Impact on conclusions.} The FlatMLP test\_acc anomaly does not affect RQ1, RQ2, or RQ3, which
use only gap Spearman values that are consistent across two independent runs. The anomaly exclusively
affects RQ4 ($\Delta$ computation), where FlatMLP test\_acc serves as the control baseline. We cannot
determine from current experiments whether the $\Delta$ null result reflects a genuine mechanism failure
or a confounded control.

\subsubsection{6.3 Limitations}

\textbf{Limitation 1: FlatMLP test\_acc substantially below literature benchmark (r = 0.279 vs. ~0.85).}
This is the most significant limitation. Our zoo does not replicate the original Unterthiner 2020
results for test\_acc prediction, limiting cross-literature comparability and confounding the $\Delta$
mechanism analysis.
\textit{Future mitigation:} Use the original Unterthiner 2020 zoo data; run 50-trial search for FlatMLP
on test\_acc; verify reproduction of r $\approx$ 0.85 before recomputing $\Delta$.

\textbf{Limitation 2: Small hyperparameter search budget (3 trials vs. 50 pre-specified).}
All Spearman values may be below encoder optimal performance. Gap results are consistent across
two independent runs (h-e1 and h-m1 agree within 5\% on FlatMLP and DWSNet), providing confidence
in the existence result. Relative rankings may shift with extended search.
\textit{Future mitigation:} 50-trial random search for all encoders on both targets; sensitivity analysis
across top-5 configurations.

\textbf{Limitation 3: Single zoo, single architecture family (Unterthiner CIFAR-10 small CNNs).}
All results are specific to this zoo. Generalization to ResNets, Transformers, or larger models is
not established.
\textit{Future mitigation:} Apply to Schürholt PDFD zoo (multi-architecture); test on ImageNet-scale zoos.

\textbf{Limitation 4: Differential advantage hypothesis not confirmed.}
The target-specificity mechanism claim is confounded. We report a null result that cannot be cleanly
interpreted as a genuine mechanism failure.
\textit{Future mitigation:} Execute corrected experiments (original data + extended search) before making
mechanism claims.

\textbf{Note on citations.} References in this paper were sourced by arXiv identifier from pipeline
artifacts. Cross-checking against Semantic Scholar or equivalent databases is recommended before
submission.

\subsubsection{6.4 Broader Impact}

Weight-based gap prediction has applications in: (1) reducing held-out test set evaluation cost
for large model zoos; (2) enabling real-time overfitting monitoring during neural architecture search
without repeated test-set evaluation; (3) providing tools for understanding where in weight space
overfitting manifests. As a potential concern, automated model selection based on predicted gap
could inadvertently select models that are overfit to the weight-space predictor's training
distribution rather than genuinely low-overfitting models. We recommend gap predictors be used
as soft filters for candidate generation, not hard selectors, in deployment settings.

